\documentclass[12pt, a4paper]{article}

% Essential Math Packages
\usepackage{amsmath}
\usepackage{amssymb}
\usepackage{amsthm}
\usepackage{geometry}
\usepackage{enumitem}

% Page Setup and Margins
\geometry{a4paper, margin=1in}

% Header and Footer Configuration
\usepackage{fancyhdr}
\pagestyle{fancy}
\fancyhf{}
\lhead{From Calculus to Cohomology}
\chead{Homework Set \# 3}
\rhead{Name: Zhu Zijian}
\lfoot{Instructor: Yuan Gao}
\rfoot{Nanjing University}
\cfoot{\thepage}

% Theorem and Problem Environments
\theoremstyle{definition}
\newtheorem{problem}{Problem}

% Custom Solution Environment (uses a black square for the QED symbol)
\newenvironment{solution}
  {\renewcommand\qedsymbol{$\blacksquare$}\begin{proof}[Solution]}
  {\end{proof}}

% Custom Commands
\newcommand{\R}{\mathbb{R}}
\newcommand{\im}{\text{im}\,}
\newcommand{\id}{\text{id}}

\begin{document}

% Academic Integrity / Collaboration Declaration
\noindent \textbf{Collaborators:} \textit{None}
\vspace{0.8cm}

% --- Problem 1 ---
\begin{problem}
\end{problem}

\begin{solution}
(a) For any $(x, y) \in \R^2$:
\begin{align*}
    (d^1 \circ d^0)(x, y) &= d^1(x, y, x - y) \\
    &= x - y - (x - y) = 0.
\end{align*}
The sequence forms a valid cochain complex.

(b) Based on \textbf{Definition 1.3}:
\begin{itemize}[nosep]
    \item $Z^0 = \ker(d^0) = \{(x, y) \mid (x, y, x-y)=(0,0,0)\} = \{0\}$.
    \item $B^0 = \im(d^{-1}) = \{0\}$.
    \item $Z^1 = \ker(d^1) = \{(u, v, w) \mid u-v-w=0\} = \{(u, v, u-v)\}$.
    \item $B^1 = \im(d^0) = \{(x, y, x-y) \mid x,y \in \R\}$. Note $Z^1 = B^1$.
    \item $Z^2 = \ker(d^2) = \R$ (since $C^3=0$).
    \item $B^2 = \im(d^1) = \R$ (since $d^1(1,0,0)=1$, the map is surjective).
\end{itemize}

(c) $H^0 = Z^0/B^0 = 0$, $H^1 = Z^1/B^1 = 0$, $H^2 = Z^2/B^2 = 0$. 
Since all cohomology groups vanish, the complex is \textbf{acyclic}.
\end{solution}

% --- Problem 2 ---
\begin{problem}
\end{problem}

\begin{solution}
$d^k_{A \oplus B}(a, b) = (d^k_A(a), d^k_B(b))$. 
\begin{align*}
    Z^k(A \oplus B) &= \{(a, b) \mid (d_A a, d_B b) = (0, 0)\} = Z^k(A) \oplus Z^k(B) \\
    B^k(A \oplus B) &= \{(d_A a', d_B b') \mid a' \in A^{k-1}, b' \in B^{k-1}\} = B^k(A) \oplus B^k(B)
\end{align*}
Then $H^k(A \oplus B) = \frac{Z^k(A) \oplus Z^k(B)}{B^k(A) \oplus B^k(B)} \cong \frac{Z^k(A)}{B^k(A)} \oplus \frac{Z^k(B)}{B^k(B)} = H^k(A) \oplus H^k(B)$.
\end{solution}

% --- Problem 3 ---
\begin{problem}
\end{problem}

\begin{solution}
(a) For any $(u, v) \in C^1$:
\begin{align*}
    (d^0 \circ h^1)(u, v) &= d^0(u) = (u, 2u) \\
    (h^2 \circ d^1)(u, v) &= h^2(2u - v) = (0, v - 2u) \\
    \implies (d^0 h^1 + h^2 d^1)(u, v) &= (u, 2u) + (0, v - 2u) = (u, v).
\end{align*}
(b) By \textbf{Lemma 1.22}, the identity map is chain homotopic to zero. Thus, $H^1(\id) = H^1(0)$, which implies the identity map on $H^1(C)$ is the zero map. Therefore, $H^1(C) = 0$.
\end{solution}

% --- Problem 4 ---
\begin{problem}
\end{problem}

\begin{solution}
(a) $d_B^0 \circ f^0(x) = 3(2x) = 6x$ and $f^1 \circ d_A^0(x) = 3(2x) = 6x$.

(b) $Cone(f)^k = A^{k+1} \oplus B^k$.
\begin{itemize}[nosep]
    \item $C^{-1} = A^0 \oplus 0 = \R$, $C^0 = A^1 \oplus B^0 = \R \oplus \R$, $C^1 = 0 \oplus B^1 = \R$.
    \item $d^{-1}(a) = (-d_A^0 a, f^0 a) = (-2a, 2a)$.
    \item $d^0(a, b) = (-d_A^1 a, f^1 a + d_B^0 b) = (0, 3a + 3b)$.
\end{itemize}

(c) $H^{-1} = \ker(d^{-1}) = 0$. $H^0 = \ker(d^0)/\im(d^{-1}) = \{(a, -a)\}/\{(-2a, 2a)\} = 0$. $H^1 = \R/\im(d^0) = \R/\R = 0$.
\end{solution}

% --- Problem 5 ---
\begin{problem}
\end{problem}

\begin{solution}
(a) $(d_C \circ i^0)(x) = d_C(x, 0) = (0, 0)$ and $(i^1 \circ d_A)(x) = i^1(0) = (0, 0)$. Stable under differential.

(b) $(C/A)^0 = \R^2/\R \cong \R$ (coord $y$), $(C/A)^1 = \R^2/\R \cong \R$ (coord $v$).
$d_{C/A}[x, y] = [d_C(x, y)] = [(y, 0)]$. Since $(y, 0) \in A^1$, its class is $[0]$. Thus $d_{C/A} = 0$.

(c) $H^0(A) = H^1(A) = \R$. $H^0(C) = \ker(d_C) = \R, H^1(C) = \R^2/\im(d_C) \cong \R$. $H^0(C/A) = H^1(C/A) = \R$.
\end{solution}

% --- Problem 6 ---
\begin{problem}
\end{problem}

\begin{solution}
The sequence does not split. As discussed in \textbf{Remark 1.32}, a splitting requires a cochain map section $s: (C/A) \to C$. 
Suppose $s^0(y) = (cy, y)$ and $s^1(v) = (dv, v)$. Commutativity requires $d_C \circ s^0 = s^1 \circ d_{C/A}$. 
\begin{align*}
    (d_C \circ s^0)(y) &= d_C(cy, y) = (y, 0) \\
    (s^1 \circ d_{C/A})(y) &= s^1(0) = (0, 0)
\end{align*}
This implies $y=0$ for all $y$, a contradiction. Thus the sequence is non-split.
\end{solution}

% --- Problem 7 ---
\begin{problem}
\end{problem}

\begin{solution}
    (a) We have $0 \to H^0(A) \to H^0(C) \to H^0(C/A) \xrightarrow{\delta^*} H^1(A) \to H^1(C) \to H^1(C/A) \to 0$. Substituting the spaces $\R$ from Problem 5(c):
    \[ 0 \to \R \xrightarrow{i^*} \R \xrightarrow{p^*} \R \xrightarrow{\delta^*} \R \xrightarrow{i^*} \R \xrightarrow{p^*} \R \to 0 \]
    (b) Following the construction in \textbf{Snake Lemma}:
    \begin{enumerate}
        \item Let $[y] \in H^0(C/A)$ be represented by $y \in (C/A)^0$.
        \item Lift $y$ to $b \in C^0$ via $p$: choose $b = (0, y)$. (Since $p(0,y) = y$).
        \item Apply $d_C$: $d_C(0, y) = (y, 0) \in C^1$.
        \item Pull back to $A^1$ via $i$: since $i^1(a) = (a, 0)$, we find $a = y$.
        \item The connecting homomorphism is defined as $\delta^*([y]) = [a] = [y] \in H^1(A)$.
    \end{enumerate}
    Since $\delta^*(y) = y$, it is the identity map $\id_\R$, which is trivially an \textbf{isomorphism}.
\end{solution}

% --- Problem 8 ---
\begin{problem}
\end{problem}

\begin{solution}
(a) $U$ is the plane minus axes, which has \textbf{4 connected components}. By \textbf{Proposition 3.1}, $\dim H^0_{dR}(U) = 4$.
(b) Basis consists of locally constant functions $f_i$ which are 1 on the $i$-th quadrant and 0 elsewhere.
\end{solution}

% --- Problem 9 ---
\begin{problem}
\end{problem}

\begin{solution}
(a) $U$ is a star-shaped area. By \textbf{Poincaré Lemma}, $H^1_{dR}(U)=0$. Thus every closed form is exact.
(b) Let $f(x, y) = \arctan(y/x)$. Then:
$\frac{\partial f}{\partial x} = \frac{-y}{x^2+y^2}$, $\frac{\partial f}{\partial y} = \frac{x}{x^2+y^2} \implies df = \omega$.
\end{solution}

% --- Problem 10 ---
\begin{problem}
\end{problem}

\begin{solution}
(a)$(\psi \circ \phi)^* = \phi^* \circ \psi^*$. Given $\psi \circ \phi = \id_U$, we have $\phi^* \circ \psi^* = (\id_U)^* = \id_{H^k_{dR}(U)}$.
(b) Since $\phi^* \circ \psi^* = \id$, $\phi^*$ has a right inverse, thus is \textbf{surjective}. $\psi^*$ has a left inverse, thus is injective.
\end{solution}

\end{document}